\section{Overview}
In this chapter, we describe  a simulation-based validation platform 
designed and implemented for self-replicating robotic systems (SRRS), 
combining the ROS 2 Jazzy 
middleware, Gazebo Harmonic physics simulation, and a modular set of 
Python controllers and sensor nodes. The principal goal of this project 
was not only to design a single robot controller, but also to create an 
experimental architecture that allows us to test and compare different control 
approaches, including  based  on NuSMV for  assembly of new robotic units.

We structured the SRRS project as a collection of ROS 2 nodes and launch 
files that collectively define the robotic world, the robot itself, 
the manipulable parts, and the control interface. The starting point of 
the simulation is the launch file \texttt{robot.launch}, which initiates the 
Gazebo world, spawns the robot and the base environment, and launches the 
graphical interface for control (robot\_controller\_gui.py). The design ensures 
that the entire system can be initialized with a single command and 
that all subsystems--control, sensing, visualization -- operate concurrently under 
ROS2 multithreaded executor.

\subsection{System Architecture }

The simulated environment consists of three categories of entities:
\begin{itemize}
    \item The main robot, defined in robot.xacro, which acts as both manipulator and replicator.
    \item Modular parts, defined in part.xacro, which represent the components that 
    the robot manipulates to assemble a new unit.
    \item The base environment, described in base.sdf, which provides a 
        fixed reference plane and a physical structure for attachment operations.
\end{itemize}

All these elements are spawned into the Gazebo Harmonic simulation using the 
ROS 2 launch mechanism.

The robot model includes multiple revolute joints 
(rev0\_1, rev2\_3, rev5\_6, rev8\_9, rev9\_10) allowing high-degree-of-
freedom manipulation. 
Each part is modeled as a single block that can 
attach to either of the robot’s two end effectors or to the base via 
virtual "attachment joints." These joints are controlled through ROS 
topics using the std\_msgs/msg/Empty message type to send attach/detach 
commands, bridged to Gazebo via ros\_gz\_bridge.

The modularity of this structure allows us to simulate replication cycles 
where the robot assembles a set of parts into a new robot instance. 
This provides a flexible testbed to evaluate different control policies-
deterministic sequences, feedback-based behaviors, or Model checking based strategies
-without modifying the core simulator.

We organized the software in several layers, from low-level motion control 
to high-level behavior orchestration:
\begin{itemize}
    \item Low-level control node: RobotController (see robotController.py)
    \item Part interface node: PartController (see partController.py)
    \item Sensor acquisition node: SRRSsensorsNode (see SRRSsensorsNode.py)
    \item Supervisory control node: SRRScontrollerNode (see SRRScontrollerNode.py)
    \item Graphical user interface: robot\_controller\_gui.py
\end{itemize}

The robot\_controller\_gui.py file launches and connects all these nodes,
forming the interactive entry point to the SRRS simulation. The GUI shown in Figure \ref{fig:simGUI} 
is implemented in Tkinter, providing manual and automated control capabilities,
live feedback of joint states, and camera streams. 

\begin{figure}[h!]
    \centering
    \includegraphics[height=8cm]{simGUI.png} 
    \caption{Graphical user interface for simulation control}
    \label{fig:simGUI}
\end{figure}

When the "Start Assemble"
button is pressed, the GUI calls start\_assemble(), which executes an 
automated sequence of motions to assemble a new robot from scattered parts.


\subsection{Robot simulation and control}

As the example of a robot we take a five-joint manipulator with two  fixed ends, 
enabling alternating attachment and detachment to the base during locomotion 
and assembly, the initial configuration is shown in Figure \ref{fig:RobotStructureGazeboSim}. 

\begin{figure}[h!]
    \centering
    \includegraphics[height=11cm]{RobotStructureGazebo.png} 
    \caption{Robot structure in Gazebo simulation environment}
    \label{fig:RobotStructureGazeboSim}
\end{figure}

This design emulates the behavior of a linear self-
reconfiguring structure similar to a modular chain robot.

We modeled the motion of the robot's free end in Cartesian coordinates 
(x, y, z) relative to the fixed base. 
The function goto\_XYZ(x, y, z) (in RobotController) defines the high-level command for moving the 
robot’s end-effector to a desired position. 
The method recursively decomposes 
the trajectory into steps of a maximum length units to prevent excessive 
joint rotations, invoking the lower-level function calculate\_joint\_angles() 
for each sub-movement.

Mathematically, for a given target point (x, y, z), we compute three internal 
angular parameters:

\[ \alpha = \arcsin(\frac{\sqrt{r_0^2+z^2}}{4}) \]
\[ \beta = \arctan(\frac{y}{x}) \]
\[ r = \sqrt{y^2+x^2} \]
\[ \gamma = \arctan(\frac{z}{r_0}) \]
,where:

$\alpha$ - is a pitch of middle block without correction 
$\beta$ - yaw of a first block
$\gamma$ - correction of $\alpha$
$ r_0=\sqrt{x^2+y^2-1} $ is a correction term depending on the 
robot geometry. 
These angles are used to derive the five joint commands:
